% !TeX root = ../main.tex
\begin{figure}[H]
\centering
\small
\textbf{Subtype constraint reduction}
\medskip

\begin{minipage}{\linewidth}
After the shape check, reduce the generated constraints below.
Resolve bound type variables before each reduction and apply substitutions
throughout the constraints.
\end{minipage}
\[
\begin{aligned}
 \tyR a<:\tyR b
   &\;\leadsto\;\{a\le b\},\\[4pt]
 (\tau_1\to\tau_2)<:(\sigma_1\to\sigma_2)
   &\;\leadsto\;\{\sigma_1<:\tau_1,\;\tau_2<:\sigma_2\},\\[4pt]
 (\tau_1\star\tau_2)<:(\sigma_1\star\sigma_2)
   &\;\leadsto\;\{\tau_1<:\sigma_1,\;\tau_2<:\sigma_2\}
       &&(\star\in\{\times,+\}),\\[4pt]
 \tyL\tau<:\tyL\sigma
   &\;\leadsto\;\{\tau<:\sigma\},\\[4pt]
 \chi<:\chi
   &\;\leadsto\;\varnothing
       &&(\chi\in\{\tyU,\tyB\}).
\end{aligned}
\]
\par\vspace{8pt}

\textbf{Constructor expansion}
\medskip

\begin{minipage}{\linewidth}
$X,Y$ are unbound type variables. For a constructor-headed type $\tau$,
let $h(\tau)$ keep only its outer constructor, with fresh children:
\end{minipage}
\[
\begin{aligned}
 h(\tyR a)&=\tyR b,
 &h(\tau\to\sigma)&=X\to Y,
 &h(\tyL\tau)&=\tyL X,\\[3pt]
 h(\tau\star\sigma)&=X\star Y\quad(\star\in\{\times,+\}),
 &h(\tyU)&=\tyU,
 &h(\tyB)&=\tyB.
\end{aligned}
\]
Variables introduced by $h$ are fresh.
\smallskip

\renewcommand{\arraystretch}{1.3}
\begin{tabular*}{\linewidth}{@{}l@{\extracolsep{\fill}}ll@{}}
\toprule
Constraint & Substitution & Remaining constraint\\
\midrule
$X<:X$ & --- & none\\
$X<:Y\quad(X\ne Y)$ & --- & defer $X<:Y$\\
$X<:\tau$ & $X\mapsto h(\tau)$ & $h(\tau)<:\tau$\\
$\tau<:X$ & $X\mapsto h(\tau)$ & $\tau<:h(\tau)$\\
\bottomrule
\end{tabular*}
\medskip

\begin{minipage}{\linewidth}
The last two rows require constructor-headed $\tau$; the two occurrences of
$h(\tau)$ in a row denote the same fresh type. Revisit deferred constraints
after each expansion.
\end{minipage}
\par\vspace{12pt}

\textbf{Affinity solving}\leanref{affinity-solving}
\medskip

\begin{minipage}{\linewidth}
With $\G\le\E$, propagate assignments until no further change:
\end{minipage}
\begin{mathpar}
 \inferrule{a\le b\\a=\E}{b:=\E}
 \and
 \inferrule{a\le b\\b=\G}{a:=\G}
\end{mathpar}
\begin{minipage}{\linewidth}
Explicit annotations and required $\G$ affinities supply the initial assignments.
Conflicting assignments reject inference. Set remaining affinity variables to
$\E$ and remaining type variables to $\tyU$.
\end{minipage}

\caption{Solving the constraints generated in \cref{fig:inference}.}
\label{fig:inference-solving}
\Description{Structural subtype constraints reduce componentwise, with
contravariance for function arguments. A preliminary shape-unification check
rejects constructor clashes and infinite types. Unbound type variables expand
to constructors with fresh children; variable inequalities remain deferred. Affinity
constraints propagate E forward and G backward; remaining affinities default
to E, and remaining type variables to unit.}
\end{figure}
